\section{Two compressed derivative calculations}
\label{sec:details}
\subsection{Eliminating the carrier envelope's vertical derivatives}
Insert $\phi_1=Ae^{k_0z}e^{i(k_0x-\omega_0t)}$ into Laplace's equation.
Differentiate both carrier factors and the slow coefficient. After the
$k_0^2$ contributions cancel, the fundamental coefficient satisfies
\begin{equation}\label{eq:vertical-envelope}
 2k_0A_z=-2ik_0A_x-(A_{xx}+A_{yy}+A_{zz}).
\end{equation}
This is (2.12). At leading slow order, $A_z\simeq-iA_x$. Therefore
$A_{zz}\simeq-A_{xx}$ and the two longitudinal second derivatives cancel.
The next iteration gives $A_z\simeq-iA_x-A_{yy}/(2k_0)$.
Because the constant-coefficient derivatives commute, it gives
\[
 A_{zz}\simeq
 \left(-i\partial_x-\frac1{2k_0}\partial_y^2\right)^2A
 =-A_{xx}+\frac i{k_0}A_{xyy}+\text{higher slow orders}.
\]
Substitution into \eqref{eq:vertical-envelope} produces (2.13):
\begin{equation}\label{eq:vertical-reduced}
 2k_0A_z=-2ik_0A_x-A_{yy}-\frac i{k_0}A_{xyy}
        +O(\varepsilon^5).
\end{equation}
Here $A$ is first order, so the retained $A_{xyy}$ term is fourth order
and omitted terms have higher total order. The paper writes its
remainder as $o(\varepsilon^4)$. The cancellation is not accidental.
A sideband's exact positive vertical exponent is
$\sqrt{(k_0+q)^2+r^2}$, which has no purely longitudinal $q^2$ term
near $k_0>0$.

\subsection{Reducing the second time derivative}
\label{sec:reduction}
In carrier-scaled variables, the left side of (2.17) is
\[
 2i\Dg A+\frac12A_{yy}-A_{tt}+\frac i2A_{xyy}.
\]
From NLS at third order,
\begin{equation}\label{eq:order-reduction-F}
 A_t=-\frac12A_x+F+O(\varepsilon^4),\qquad
 F=\frac i4A_{yy}-\frac i8A_{xx}-\frac i2|A|^2A.
\end{equation}
$F$ is third order. At the order needed to differentiate it once,
the envelope translates at velocity $1/2$, so $F_t=-F_x/2$ up to
fifth-order terms. Thus
\begin{align*}
 A_{tt}&=-\frac12A_{xt}+F_t+O(\varepsilon^5)\\
 &=\frac14A_{xx}-\frac12F_x+F_t+O(\varepsilon^5)\\
 &=\frac14A_{xx}-\frac i4A_{xyy}+\frac i8A_{xxx}
   +\frac i2\partial_x(|A|^2A)+O(\varepsilon^5).
\end{align*}
This yields the linear correction
$-i(6A_{xyy}-A_{xxx})/8$ on the right of the reduced NLS equation.
The derivative of the cubic term also changes nonlinear gradient coefficients.

To expose the source-consistency issue without concealing it, write the
dimensionless right-side coefficient of the intermediate equation as
\[
 \Gamma=|A|^2+i\beta(AA_x^*-A^*A_x)
              -i\partial_x|A|^2+\Ph_x-i\Ph_z.
\]
The $3ik_0^3$ coefficient printed below (2.17), with the scaling (2.18),
gives $\beta=3/4$. After reducing $A_{tt}$, the nonlinear gradient terms
on the right become
\begin{equation}\label{eq:beta-reduction}
 i(\beta-\tfrac12)A(AA_x^*-A^*A_x)
 -\frac i2|A|^2A_x.
\end{equation}
For $\beta=3/4$, the first coefficient is $i/4$.
The scan's (2.19) instead prints $3i/2$ and $+i/2$.
To obtain the linearized coefficients in (3.2)--(3.3),
\eqref{eq:beta-reduction} would require $\beta=2$ and its displayed minus
sign. We have not established an authoritative corrected edition or erratum;
these are direct algebraic comparisons within the supplied scan.

The practical lesson is to distinguish a formal method you can learn from
a printed coefficient chain that may need source correction. The stability
determinant derived in Section~\ref{sec:stability} is independently checked
from the stated (3.2)--(3.3). The mean-flow mechanism and NLS prerequisite
calculations do not depend on resolving this coefficient issue.

\subsection{Verify the discrepancy by linearizing one gradient term}
Let the right side of the full NLS balance contain
\[
 ic\,A(AA_x^*-A^*A_x)+id\,|A|^2A_x+A(\Ph_x-i\Ph_z).
\]
For $A=A_0e^{-iA_0^2t/2}(1+\eta u+i\eta v)$, its linear gradient part
is proportional to
\[
 A_0^3\eta\bigl[id\,u_x+(2c-d)v_x\bigr].
\]
Separating the envelope equation into its real and imaginary parts gives
\[
 P_1u-\frac d2A_0^2u_x+P_2v=-\frac12\Ph_z,
 \qquad
 P_1v+(c-\tfrac d2)A_0^2v_x-P_2u+A_0^2u=-\frac12\Ph_x.
\]
The printed $c=3/2,d=+1/2$ yield $-1/4,5/4$ before eliminating
$\Ph_z$. The paper's (3.2)--(3.3) use $+1/4,7/4$, corresponding to
$c=3/2,d=-1/2$. This supplies a short reproducible check of the sign issue.
\clearpage
